\section{Wick products on larger cutoff tori}\label{sec:vtm-wick}

Use the cell sources of Section~\ref{sec:reference-cells}
and write
\[
 N_*=8L^3,\qquad
 a_L=\frac{L^3-1}{(L-1)^3},\qquad
 D^{\mathrm{mat}}=L^3\log z_{r,n_*}-\log z_{r,Ln_*},
 \qquad
 \delta=\frac{D_L}{(L-1)^3N_*}.
\]
Fix $0<e\le1/3$, with $q=e\sqrt\eta$ and
$q^2n_*=e^2R^*$, and assume
\begin{equation}\label{eq:transfer-W3}
 6N_*e^2R^*\log(1/e)\le\pi^2.
\end{equation}
The choice of scales ensuring this condition is given
below. Let $\mathcal C_*=\{0,\ldots,2L-1\}^3$
and use the probe variance $v_\eta$ of
\eqref{eq:cells-orthogonality}. For real $c_2,c_4$ put
\begin{align}
 O_m^{[c]}&=W_m(X_c;v_\eta,c_m),\qquad m=2,4,
                                      \label{eq:transfer-Wick-observable}\\
 \Theta_m^{(0)}
  &=\left|
      \left[\prod_{c\in\mathcal C_*}
            (\partial_{t_c}^m-c_m)\right]\mathring\Phi(0)
    \right|^{1/N_*}.                  \label{eq:transfer-Theta}\\
 M_4&=600N_*^2v_\eta^2e^{-4}+|c_4|,\qquad
 M_2=12N_*v_\eta e^{-2}+|c_2|,\notag\\
 T_m&=\Theta_m^{(0)}+2e^6M_m,\qquad
 u_m=\exp(\delta)T_m.                 \label{eq:transfer-Wick-data}
\end{align}
Here $\mathring\Phi$ is the normalized external partition
function \eqref{eq:external-reference-partition}.
The notation $e^{-m}$ denotes an inverse power of the
charge.

\begin{proposition}\label{prop:transfer-Wick}
Assume the probe placement condition \eqref{eq:cells-fit}.
For $n=n_*$, or $n\in p_{\rm cell}\Z$ with $n\ge Ln_*$,
the all-flux law with the fixed real lattice mass obeys
\begin{align}
 0\le\delta
 &\le\frac{a_L}{N_*}\log\frac{2\pi}{e^2R^*}
       +\frac{K_L+D^{\mathrm{mat}}+1}{(L-1)^3N_*},
                                           \label{eq:transfer-delta}\\
 \left|\E_n^{\mathrm{af}}\prod_{c\in\mathcal C}
                        O_m^{[c]}\right|
 &\le u_m^{|\mathcal C|}              \label{eq:transfer-Wick-product}
\end{align}
for every set $\mathcal C$ of distinct cells and $m=2,4$.
In particular, if $\kappa_j$ denotes the cumulant of
$X_0$ on that torus, then
\begin{align}
 \left|e^{-2}(\kappa_2-v_\eta)-c_2\right|&\le u_2,
                                      \label{eq:transfer-second-cumulant}\\
 \left|e^{-4}\kappa_4-(c_4-3c_2^2)\right|
   &\le u_4+3u_2(2|c_2|+u_2).        \label{eq:transfer-fourth-cumulant}
\end{align}
The bounds are uniform in $n$.
Their right sides are formed on the two reference tori;
$\Theta_m^{(0)}$, and therefore $u_m$, may depend on
the placement.
\end{proposition}

\begin{proof}
Condition \eqref{eq:transfer-W3} implies $e^2R^*<1$,
since $N_*\ge1000$ and $\log(1/e)\ge\log3>1$.
Put $x=\exp(-2\pi^2/(e^2R^*))$. The same condition gives
\[
 x\le e^{12N_*},\qquad x^{1/(2N_*)}\le e^6.
\]
Apply Proposition~\ref{prop:transfer-flux-defect} at
$n_*$ and $j=L$. Since $\varsigma_{n_*}\le7x<1/8$,
\[
 D_L^{\mathrm{fl}}\le-L^3\log(1-\varsigma_{n_*})
              \le8L^3x=N_*x
              \le N_*3^{-12N_*}\le1.
\]
The photon estimate and the defect decomposition
now prove the upper bound in \eqref{eq:transfer-delta};
its lower bound is Proposition~\ref{prop:transfer-general}.

All the polynomials in \eqref{eq:transfer-Wick-observable}
are the same local integrands on every carrying torus,
by Section~\ref{sec:reference-cells}. They are real
elements of the cell observable class. Write
$O_m^\sharp=\prod_{c\in\mathcal C_*}O_m^{[c]}$ on
$C(n_*)$. Its all-flux expectation is nonnegative,
by the fold trace \eqref{eq:transfer-fold-trace}.
We estimate it by splitting flux zero from the
other sectors.

First, Gaussian domination and the moment bound
give, for every positive integer $i$,
\[
 \E_{n_*}^{\mathrm{af}}X_c^{2i}
                              \le(2i\exp(1)v_\eta)^i.
\]
Minkowski's inequality at exponent $2N_*$ yields
\begin{align*}
 \|H_4(X_c;v_\eta)\|_{2N_*}
   &\le[64\exp(2)N_*^2+24\exp(1)N_*+3]v_\eta^2
       \le600N_*^2v_\eta^2,\\
 \|H_2(X_c;v_\eta)\|_{2N_*}
   &\le[4\exp(1)N_*+1]v_\eta
       \le12N_*v_\eta.
\end{align*}
Thus $\|O_m^{[c]}\|_{2N_*}\le M_m$.
The finite moments justify all uses of H\"older
and Cauchy--Schwarz below, also when $v_\eta=0$.

On the zero-flux sector,
\eqref{eq:external-Wick-product} gives
\[
 \left|\E_{n_*}^{\mathrm{af}}
               [O_m^\sharp\mathbf1_{\{\mathrm{flux}=0\}}]\right|
   =P_{n_*}(0)\left|\E_{n_*}^{(0)}O_m^\sharp\right|
   \le(\Theta_m^{(0)})^{N_*}.
\]
On the remaining sectors, Cauchy--Schwarz followed
by H\"older with $N_*$ factors gives
\begin{align*}
 \left|\E_{n_*}^{\mathrm{af}}
               [O_m^\sharp\mathbf1_{\{\mathrm{flux}\ne0\}}]\right|
 &\le\varsigma_{n_*}^{1/2}
                \bigl(\E_{n_*}^{\mathrm{af}}|O_m^\sharp|^2\bigr)^{1/2}\\
 &\le\varsigma_{n_*}^{1/2}
                \prod_{c\in\mathcal C_*}\|O_m^{[c]}\|_{2N_*}\\
 &\le\bigl(\varsigma_{n_*}^{1/(2N_*)}M_m\bigr)^{N_*}.
\end{align*}
Since $\varsigma_{n_*}^{1/(2N_*)}\le2e^6$,
the nonnegative $N_*$-th root satisfies
\[
 \|O_m\|_*=
       \bigl(\E_{n_*}^{\mathrm{af}}O_m^\sharp\bigr)^{1/N_*}
             \le\Theta_m^{(0)}+2e^6M_m=T_m.
\]
For nonnegative $a,b$ we used
$(a^{N_*}+b^{N_*})^{1/N_*}\le a+b$.
Proposition~\ref{prop:transfer-general}, with
$s_0=2L$ and $j=L$, proves
\eqref{eq:transfer-Wick-product}.

For one cell the same estimate gives
$|e^{-m}\E H_m(X_0;v_\eta)-c_m|\le u_m$.
Charge conjugation makes the odd moments zero.
Writing $y=e^{-2}\E H_2(X_0;v_\eta)$, one has
\[
 e^{-4}\kappa_4=e^{-4}\E H_4(X_0;v_\eta)-3y^2.
\]
The second-order conclusion follows at once.
For the fourth, use $|y-c_2|\le u_2$ and
$|y+c_2|\le2|c_2|+u_2$ to obtain
\eqref{eq:transfer-fourth-cumulant}.
\end{proof}

